\section{Benchmark results}
\label{sec:benchmark}

Across the 600 item-model prompts of \GeoGenBench, the structured mode raised automatic pass rates for every model that could run both modes (\cref{tab:results}). \GeoGenBench holds 801 prompts, all scored by the checker of \cref{sec:scoring}: 600 generated from the item models described above, and 201 written from the topic scaffold of a K--12 mathematics curriculum~\citep{bluebonnet2024}. The results below cover the 600. The prompt set, item models, checker vocabulary, and rendered outputs are public, and so is the code.\footnote{Data: \url{https://huggingface.co/datasets/meix/geogenbench}. Code: the repository in footnote 1.} We ran six large models from three companies, each in two modes. The \emph{by-hand} mode has the language model write drawing code directly. The \emph{structured} mode has it fill in the description of \cref{sec:pipeline}. Each prompt was run three or more times, so a cell in \cref{tab:results} holds 1,800 to 2,861 generations, and a rate is the share of those generations that passed the automatic checker. The checker overestimated human acceptance in the validation sample (\cref{sec:scoring}), so these rates may overestimate correctness.

\paragraph{How we ran it.}
The by-hand mode asks the language model for TikZ drawing code. The structured mode asks for the typed form of \cref{sec:pipeline}. For the by-hand mode, an extractor recovers the coordinates from the TikZ source; output the extractor cannot parse counts as a failure. Both run at temperature $1.0$, the sampling setting that controls randomness, with no explicit limit on output length. Both have the same budget of up to three attempts. What differs is the retry trigger. The by-hand mode retries only when its code fails to render. The structured mode also retries when its description fails to compile or the figure breaks a check it declared. Those declared checks are the language model's own. The held-out scoring predicates generated by the item model, the ones we grade against, are never shown to the language model. The comparison is therefore between two complete workflows. The structured mode's feedback lets a model catch its own geometric mistakes, which the by-hand mode cannot, and its extra retries appear in its per-item cost. Each requirement is decided on the computed coordinates within an absolute tolerance of $\tau = 5\times10^{-3}$, applied to angles in radians and to lengths in coordinates normalised to $[-10, 10]$. The $600$ item-model prompts divide evenly across three difficulty tiers of $200$ each. They come from $30$ item models covering triangles, quadrilaterals, circles, segments, and lines, with $3.2$ requirements apiece on average. The $201$ curriculum prompts were drafted with a language model, with answer keys written by hand. They carry more requirements, $6.0$ on average. Cost per item is the input and output tokens of the prompt and its response, summed across any retries and priced at each vendor's published rates, then averaged over every prompt, failures included. Failures are dominated by wrong geometry, a requirement that does not hold, rather than by broken drawing code. The structured mode's remaining failures are mostly timeouts.

\begin{table}[t]
\centering
\footnotesize
\setlength{\tabcolsep}{5pt}
\begin{tabular}{@{}lrrr@{}}
\toprule
\textbf{Model} & \textbf{By hand} & \textbf{Structured} & \textbf{\$/item}\\
\midrule
GPT-5.5            & 91.9\% & \textbf{94.1\%} & 0.075\\
Claude Sonnet 4.6  & 81.2\% & 93.4\% & 0.088\\
Claude Opus 4.7    & 71.6\% & 93.3\% & 0.53\\
Claude Haiku 4.5   & 70.5\% & 75.3\% & 0.032\\
Gemini 2.5 Pro     & 58.4\% & 83.1\% & 0.050\\
Gemini 2.5 Flash   & 35.9\% & blocked & ---\\
\midrule
\emph{Confirmed by graders} & 75\% & 87\% & ---\\
\bottomrule
\end{tabular}
\caption{Share of generations \emph{passing the current automatic checker} over the $600$ item-model prompts, each run three or more times, in the by-hand mode and the structured mode, with the average cost of one structured generation in US dollars. The checker overestimated human acceptance in the validation sample (\cref{sec:scoring}), so these rates may overestimate correctness. The last row gives the share of automatic passes that human graders confirmed, pooled over models. Every model that could run the structured mode improves. The three strongest converge near $93$--$94\%$ despite a sevenfold difference in cost. Gemini~2.5~Flash could not run the structured mode (see text). Bold marks the highest rate.}
\label{tab:results}
\end{table}

The structured mode helped every model that could run it (\cref{tab:results}). Pass rates rose by 2 to 25 percentage points. The largest gain, 25 points, went to Gemini~2.5~Pro, the weakest by-hand drawer among the models that could run both modes. In the structured mode, the three strongest models finished within a single point of one another, despite a sevenfold difference in price.

The human study of \cref{sec:scoring} supports a sensitivity calculation. In that study the checker erred in one direction only. The graders rated diagrams from a pilot run of the same benchmark. They confirmed 75\% of the automatic passes drawn by hand (86 of 114) and 87\% of those from the structured mode (55 of 63). If those proportions carry over to the full run, the structured rates in \cref{tab:results} fall to about seven-eighths of those shown and the by-hand rates to about three-quarters. Under that assumption the structured advantage survives and widens, which would mean the structured mode also reduces errors the checker misses. The per-model cells in that study hold 7 to 23 diagrams, too few to correct each row separately.

One practical limitation concerns vendor infrastructure. Gemini~2.5~Flash could not run the structured mode. The API rejected our output schema, with its 31 kinds of building block, as too complex before generation began, with the message that the schema ``produces a constraint that has too many states for serving.'' The limit is in the vendor's constrained-decoding service, not in the model. Gemini~2.5~Pro and GPT-5.5 accepted the same schema. The results support a single point: the structured mode helps at scale, across every model that could run it.
